\section*{ЗАКЛЮЧЕНИЕ}
\addcontentsline{toc}{section}{ЗАКЛЮЧЕНИЕ}

В ходе прохождения практики было разработано программное обеспечение для оценки кредитоспособности физических лиц с использованием алгоритма градиентного бустинга деревьев решений. Разработанный программный комплекс включает пользовательский интерфейс, реализованный на языке C\# с использованием технологии Windows Forms, и модуль предсказания, реализованный на языке Python.

В рамках работы была реализована возможность ручного ввода данных заемщика, загрузки данных из JSON-файла, формирования входной структуры данных, запуска Python-модуля предсказания, получения результата скоринга и отображения итоговой оценки пользователю. Результат работы программы включает вероятность дефолта, скоринговый балл, уровень кредитного риска, факторы риска и рекомендации.

В рамках работы была проведена проверка функциональных возможностей разработанного программного обеспечения. В процессе тестирования проверялись основные сценарии работы приложения: ввод сведений о заемщике через формы приложения, создание JSON-структуры, загрузка данных из файла, запуск Python-скрипта из приложения, получение результата скоринга, отображение результата пользователю, а также обработка некорректных входных данных. По результатам функционального тестирования установлено, что разработанное программное обеспечение корректно выполняет основные пользовательские сценарии и не завершается аварийно при возникновении ошибочных ситуаций.

Для проверки входных данных были выделены классы эквивалентности, включающие корректные данные заемщика, незаполненные обязательные поля, некорректный формат данных, корректный и некорректный JSON-файл, а также различные профили заемщиков с низким, средним и высоким уровнем риска. Использование таких классов позволило проверить основные варианты поведения системы без избыточного увеличения количества тестовых случаев.

Модульное тестирование в рамках практики не проводилось. Тестирование графического интерфейса также не проводилось.

Также было проведено исследование характеристик разработанного программного обеспечения. В рамках исследования оценивалась способность модели выявлять заемщиков с высоким риском дефолта, анализировалось влияние порога классификации на количество ошибок, а также рассматривалась применимость разработанного программного комплекса для предварительной оценки кредитного риска.

По результатам исследования итоговая модель показала высокую полноту выявления дефолтных заемщиков. Значение Recall составило 0.89, что свидетельствует о способности модели выявлять значительную часть заемщиков с высоким риском дефолта. Значение ROC-AUC составило 0.777. Это показывает, что модель ранжирует заемщиков по уровню кредитного риска лучше случайного классификатора.

В ходе анализа было выявлено, что выбранный порог классификации, равный 0.3, позволяет уменьшить количество пропущенных дефолтных заемщиков. При этом увеличение порога классификации снижает количество ложноположительных решений, но одновременно увеличивает число пропущенных дефолтов. Для задачи кредитного скоринга выбранный подход является обоснованным, поскольку пропуск заемщика с высоким риском дефолта может привести к финансовым потерям.

Таким образом, результаты тестирования подтвердили корректность работы основных компонентов разработанного программного обеспечения, а результаты исследования показали применимость выбранного подхода для предварительной оценки кредитного риска физических лиц.